% Generated draft: requires mathematical and source review.
\documentclass[11pt]{article}
\usepackage[margin=25mm]{geometry}
\usepackage{fontspec}
\setmainfont{Latin Modern Roman}
\usepackage{amsmath,amssymb,mathtools,unicode-math}
\setmathfont{Latin Modern Math}
\usepackage{enumitem,needspace,xurl,hyperref}
\setlength{\parindent}{0pt}
\setlength{\parskip}{5pt}
\setlength{\emergencystretch}{3em}
\newcommand{\N}{\mathbb{N}}
\newcommand{\Z}{\mathbb{Z}}
\newcommand{\Q}{\mathbb{Q}}
\newcommand{\R}{\mathbb{R}}
\newcommand{\C}{\mathbb{C}}
\newcommand{\F}{\mathbb{F}}
\newcommand{\A}{\mathbb{A}}
\newcommand{\PP}{\mathbb{P}}
\newcommand{\E}{\mathbb{E}}
\newcommand{\eps}{\varepsilon}
\newcommand{\dd}{\,\mathrm{d}}
\begin{document}
\section*{UnsolvedMath ID: 800486}
\textbf{Source problem code:} catalogue entry
\subsection{Definitions and mathematical statement}
Let \(\gamma\subset\R^2\) be an oval enclosing a compact convex region \(K\) with nonempty interior. The source discusses smooth convex closed curves, but does not specify the precise regularity or strict convexity required of the oval. For measurable planar regions \(A,B\), define their symmetric difference and its area by
\[
A\mathbin{\triangle}B=(A\setminus B)\cup(B\setminus A),\qquad d_{\triangle}(A,B)=\operatorname{area}(A\mathbin{\triangle}B).
\]
For an integer \(n\geq3\), let \(\mathcal P_n\) denote a chosen nonempty class of admissible filled \(n\)-gons. A variational formulation of the approximation function is
\[
F_{K,\mathcal P}(n)=\inf_{P\in\mathcal P_n}\operatorname{area}(K\mathbin{\triangle}P).
\]
This expression is conditional on specifying \(\mathcal P_n\): the source does not explicitly require convex polygons, prescribe an inscribed or circumscribed constraint, or state whether exactly \(n\) or at most \(n\) vertices are allowed. No attainment of the infimum is assumed.

\textbf{Problem (Serge Tabachnikov, Problem 8).} Is the resulting optimal symmetric-difference approximation function of \(n\) related to a billiard-like dynamical system associated with the oval? The source does not define either the permitted class of dynamical systems or the intended meaning of ``related,'' so it does not specify a precise identity or existence assertion.

The supplied motivation concerns four other approximation functions, obtained by combining perimeter or area discrepancy with inscribed or circumscribed polygons. The source reports that these correspond to minimal actions of conventional (Birkhoff) billiards, outer billiards, symplectic billiards, and outer length billiards. This motivates the question but does not specify an analogous action for symmetric-difference approximation.
\subsection{Short English statement}
For a convex oval, minimize the area of the symmetric difference between its enclosed region and an approximating \(n\)-gon. Serge Tabachnikov asks whether this function of \(n\) has a billiard-like dynamical interpretation. The admissible polygons, the oval's precise regularity, and the intended dynamical relationship remain unspecified.
\subsection{Sources}
\begin{itemize}
\item Local UnsolvedMath registry, permanent ID \(800486\), source problem code \texttt{catalogue entry}, in the AMR Open Problem Lists collection. The supplied record labels the problem \texttt{open}; this is a supplied status, not an independently verified literature assessment.
\item Albach, Barraud et al., \emph{Open Problems in Billiards and Quantitative Symplectic Geometry} (2026), Serge Tabachnikov, Problem 8. Supplied source URL: \url{https://arxiv.org/abs/2602.12896}. The record attributes the contribution to Serge Tabachnikov and lists Albach, Barraud et al. as proposers. Supplied access date: 2026-10-16; extraction method: \texttt{arxiv\_latex}. These are provenance supplied with the record, not a claim of independent source access or verification.
\item Reference label [Gr] is cited in the supplied statement for approximation of smooth convex closed curves by polygons. No bibliographic details or separate link are supplied.
\item Reference labels [AT24] and [BBF] are cited in the supplied statement for the lesser studied billiard systems. No bibliographic details or separate links are supplied.
\end{itemize}
\end{document}
